\documentclass[conference]{IEEEtran}

\usepackage[T1]{fontenc}
\usepackage[utf8]{inputenc}
\usepackage{cite}
\usepackage{amsmath,amssymb,amsfonts}
\usepackage{graphicx}
\usepackage{booktabs}
\usepackage{multirow}
\usepackage{array}
\usepackage{tabularx}
\usepackage{siunitx}
\usepackage{url}
\usepackage{xcolor}
\usepackage{balance}

% Compact TODO macro for tables
\newcommand{\tbd}{---}

% Table typography macros
\newcommand{\tablesize}{\footnotesize}
\newcommand{\tablecolsep}{3pt}

\title{Bounded Residual Reinforcement Learning over Limited LQR/IK Prior for Height-Adaptive Wheeled Biped Balance}

\author{
\IEEEauthorblockN{Van Thuong Nguyen}
\IEEEauthorblockA{
Independent Researcher\\
Email: vanthuong180702t@gmail.com
}
}

\begin{document}

\maketitle

\begin{abstract}
Wheeled bipedal robots combine the efficiency of wheeled locomotion with the adaptability of articulated legs, but maintaining balance across varying body heights remains challenging.
Classical model-based controllers such as LQR with height-dependent inverse kinematics provide limited structured action priors, but are insufficient as standalone solutions---our evaluation shows fall rates of 85--100\% across the feasible height range (0.40--0.70\,m).
Pure reinforcement learning from scratch can learn robust policies but suffers from poor sample efficiency and lacks the structure that domain knowledge provides.
We propose a bounded residual reinforcement learning approach that combines a height-dependent LQR/IK prior with a learned PPO residual policy.
The prior provides limited structured action priors, while the bounded residual learns corrective actions for height-adaptive stabilization, commanded height transitions, and push-disturbance recovery.
The residual policy receives a 52-dimensional observation that includes the 42-dimensional base observation plus the 10-dimensional base action from the prior, and outputs a bounded residual action that is composed with the base action via $\text{final} = \text{clip}(\text{base} + \text{scale} \times \text{residual}, -1, 1)$.
We train the policy in MuJoCo MJX with JAX-accelerated vectorized simulation using proximal policy optimization across 4,096 parallel environments.
Our evaluation of the LQR/IK prior alone shows fall rates of 85--100\% across the feasible height range, confirming its role as a limited structured prior.
% [TODO-RESULT: Once residual training completes, add evaluation results for random-height balance, height transitions, push recovery, and robustness scenarios]
The proposed approach will be evaluated against the LQR/IK prior alone and a pure PPO baseline trained from scratch.
\end{abstract}

\begin{IEEEkeywords}
Wheeled biped robot, residual reinforcement learning, bounded residual policy, LQR/IK prior,
height-adaptive balance, push recovery, MuJoCo MJX.
\end{IEEEkeywords}

% ============================================================
\section{Introduction}
% ============================================================
Wheeled biped robots combine the efficiency of wheel-based locomotion with
the adaptability of legged morphology.
Compared with conventional legged systems, wheeled bipeds can achieve
lower-cost motion on flat terrain while exploiting articulated legs to modify
body posture, regulate center-of-mass height, and improve disturbance rejection.
However, these benefits come at the cost of more challenging control:
balancing at different commanded heights requires managing the nonlinear
coupling between wheel dynamics, torso attitude, and multi-joint leg motions.

Classical control approaches~\cite{klemm2020lqrwbc} can stabilize wheeled bipeds
near a nominal operating point, but become increasingly inaccurate as the
robot moves away from the linearization height.
Recent progress in reinforcement learning (RL) has shown that robust control
policies can be learned directly in simulation for complex robotic
systems~\cite{hwangbo2019learning,lee2020learning,rudin2022walk}.
Nevertheless, training a single policy to balance across a wide range of
body heights is difficult due to poor sample efficiency and the risk of
reward exploitation.
In particular, simple height-tracking rewards may allow policies to satisfy
torso height objectives through unnatural or unstable joint configurations
rather than through physically meaningful postures.

To address these challenges, we propose a bounded residual reinforcement learning
approach that combines a height-dependent LQR/IK prior with a learned PPO residual policy.
The prior provides limited structured action priors based on linearized dynamics
and inverse kinematics, while the bounded residual learns corrective actions to
compensate for modeling errors, handle height transitions, and recover from push disturbances.
Our evaluation of the LQR/IK prior alone shows fall rates of 85--100\% across the
feasible height range (0.40--0.70\,m), confirming that the prior is insufficient as
a standalone controller but provides useful structure for residual learning.
The residual policy receives a 52-dimensional observation that includes the base
42-dimensional proprioceptive state plus the 10-dimensional base action from the prior,
enabling the policy to learn corrections that are aware of the nominal controller's intent.

The main contributions of this work are:
\begin{itemize}
    \item A bounded residual reinforcement learning architecture that combines
    a height-dependent LQR/IK prior with a learned PPO residual policy for
    height-adaptive wheeled biped balance.
    \item Quantitative evaluation showing that the LQR/IK prior alone achieves
    85--100\% fall rates across the feasible height range (0.40--0.70\,m),
    establishing it as a limited structured prior suitable
    for residual learning.
    \item A 52-dimensional residual observation space that includes the base
    42-dimensional proprioceptive state plus the 10-dimensional base action
    from the prior, enabling the residual policy to learn corrections that
    are aware of the nominal controller's intent.
    \item Bounded residual action composition via
    $\text{final} = \text{clip}(\text{base} + \text{scale} \times \text{residual}, -1, 1)$
    with per-joint residual scale vectors, preventing the residual from
    dominating or destabilizing the base controller.
    \item Evaluation protocol for random-height balance, commanded
    height transitions, push recovery up to 200\,N, and robustness to model
    uncertainty, with comparison against the LQR/IK prior alone and a pure PPO baseline.
    % [TODO-RESULT: Add quantitative summary once residual training completes]
\end{itemize}

The rest of this paper is organized as follows.
Section~II reviews related work.
Section~III presents the robot platform and simulation setup.
Section~IV describes the bounded residual RL architecture.
Section~V details the height-dependent LQR/IK prior.
Section~VI presents the residual policy design and training.
Section~VII describes the reward formulation.
Section~VIII details the evaluation protocol and metrics.
Section~IX presents experimental results.
Section~X discusses limitations and future work.
Section~XI concludes the paper.

% ============================================================
\section{Related Work}
\label{sec:related}
% ============================================================

\subsection{Classical and Model-Based Control for Wheeled Bipeds}
Wheeled bipedal robots combine the energy efficiency of wheel-based locomotion
with the postural adaptability of articulated legs.
Early and influential work by Klemm~\emph{et al.}~\cite{klemm2019ascento}
introduced Ascento, a two-wheeled jumping robot demonstrating active balance
through LQR-based stabilization around an inverted pendulum model.
Subsequent work~\cite{klemm2020lqrwbc} extended this to an LQR-assisted
whole-body controller with kinematic loops.
Hsu~\emph{et al.}~\cite{hsu2022wbrfuzzy} applied fuzzy logic control to
wheeled biped balance regulation.
Li~\emph{et al.}~\cite{li2021wheelbiped} investigated dynamic motion control
for wheel-biped transformable robots.
The DIABLO platform~\cite{diablo2024} demonstrated six-DOF direct-drive
wheeled biped control using cascaded LQR.
While these model-based approaches provide stability guarantees near an
operating point, they require accurate dynamics models and degrade when the
robot deviates significantly from the linearization configuration---a
limitation that becomes severe under variable-height commands.

\subsection{Deep Reinforcement Learning for Legged Locomotion}
Deep RL has produced robust locomotion policies for complex legged systems.
Hwangbo~\emph{et al.}~\cite{hwangbo2019learning} introduced a learned
actuator network that accurately models motor dynamics, enabling sim-to-real
transfer of agile skills on ANYmal.
Lee~\emph{et al.}~\cite{lee2020learning} demonstrated blind proprioceptive
locomotion over challenging outdoor terrain using a teacher-student
distillation pipeline.
Kumar~\emph{et al.}~\cite{kumar2021rma} proposed Rapid Motor Adaptation (RMA),
a two-stage approach that trains an adaptation module online to handle
unseen dynamics such as varying payload and slippery floors.
Rudin~\emph{et al.}~\cite{rudin2022walk} showed that massively parallel
GPU-based simulation (4,096 environments in Isaac Gym) can train walking
policies in minutes using a terrain curriculum closely related to ours.
Miki~\emph{et al.}~\cite{miki2022wild} further demonstrated perception-aware
locomotion over kilometer-scale outdoor terrain using attention-based
proprioceptive encoders.
These works establish that RL can scale to high-dimensional contact-rich
control problems; our work adapts this paradigm to the specific challenge of
variable-height wheeled-biped balancing.

\subsection{Residual Reinforcement Learning}
Residual RL combines model-based controllers with learned policies to leverage
the benefits of both approaches.
Silver~\emph{et al.}~\cite{silver2018residual} introduced residual policy learning,
where a learned policy outputs corrections to a fixed base controller.
Johannink~\emph{et al.}~\cite{johannink2019residual} demonstrated residual RL
for robotic manipulation, showing that learning on top of a nominal controller
improves sample efficiency and final performance compared to learning from scratch.
Rana~\emph{et al.}~\cite{rana2020bayesianresidual} proposed Bayesian residual RL
for safe exploration by maintaining uncertainty estimates over the residual policy.
Our work follows this paradigm: we train a bounded PPO residual policy over a
height-dependent LQR/IK prior, with explicit action composition
$\text{final} = \text{clip}(\text{base} + \text{scale} \times \text{residual}, -1, 1)$
that prevents the residual from dominating or destabilizing the base controller.
Unlike prior work that assumes a strong base controller, we quantitatively
establish that our LQR/IK prior alone achieves 85--100\% fall rates across
the feasible height range, confirming it as a limited but structured nominal
controller suitable for residual learning.

\subsection{Reinforcement Learning for Wheeled-Legged Robots}
More recently, RL has been applied to wheeled-legged hybrid platforms.
Lee~\emph{et al.}~\cite{lee2024robustnavigation} presented a full autonomy
stack for ANYmal-on-wheels, achieving 20\,km/h driving with 83\% reduced
cost-of-transport.
Vollenweider~\emph{et al.}~\cite{vollenweider2023advancedskills} trained diverse
locomotion skills for wheeled-legged robots through adversarial motion priors.
Chamorro~\emph{et al.}~\cite{chamorro2024stairclimbing} used asymmetric
actor-critic RL with privileged terrain information for blind stair climbing
on Ascento, successfully transferring to hardware for 15\,cm step negotiation.
Zhu~\emph{et al.}~\cite{whleaper2024} introduced Whleaper, a 10-DOF wheeled
biped most similar to our platform, demonstrating both LQR and PPO control
for multimodal locomotion.
Zhu and Hayashibe~\cite{zhu2022hierarchical} applied hierarchical RL for
path tracking on a wheeled biped with a sim-to-real transfer pipeline.
In contrast to these platforms---which either fix body height or focus on
forward locomotion---our work specifically addresses \emph{variable-height
stationary balance}, requiring the policy to operate stably across a 30\,cm
height range while maintaining postural integrity.

\subsection{Curriculum Learning for Robot Control}
Curriculum learning, introduced by Bengio~\emph{et al.}~\cite{bengio2009curriculum},
organizes training from easy to hard examples to improve convergence
and generalization.
In robotics, Rudin~\emph{et al.}~\cite{rudin2022walk} demonstrated a
terrain difficulty curriculum that substantially improves traversal capability.
The Automatic Curriculum Learning (ACL) survey by
Portelas~\emph{et al.}~\cite{portelas2020acl} systematizes curriculum strategies
by learning progress, distinguishing goal-parameterized curricula from
task-space curricula.
Teacher algorithms such as ALP-GMM~\cite{portelas2020teacher} generate
parameterized task distributions that maximize learning progress in
continuous task spaces.
POET~\cite{wang2019poet} co-evolves environments and agents to produce
increasingly complex curricula.
Our within-stage height curriculum is a monotone threshold-gated curriculum:
the commanded height range expands only when a greedy evaluation pass
confirms the policy has mastered the current range, directly addressing
the instability that arises when a single policy must balance across
configurations with qualitatively different inverted-pendulum dynamics.

\subsection{Domain Randomization and Sim-to-Real Transfer}
Sim-to-real transfer for locomotion policies relies on domain randomization
to cover the reality gap.
Tobin~\emph{et al.}~\cite{tobin2017domain} established domain randomization
as a viable transfer strategy.
Peng~\emph{et al.}~\cite{peng2018simtoreal} demonstrated dynamics
randomization for robust quadruped locomotion on physical hardware.
Tan~\emph{et al.}~\cite{tan2018simtoreal} showed that randomizing actuator
parameters and adding observation noise enables Minitaur to transfer
sim-trained trotting gaits to the real robot without fine-tuning.
Our framework adopts these established protocols~\cite{tobin2017domain,peng2018simtoreal,tan2018simtoreal}---randomizing friction,
mass properties, joint damping, and adding IMU/encoder noise and
optional action delay---to prepare the trained policy for future hardware deployment.

% ============================================================
\section{Robot Platform and Simulation Setup}
% ============================================================

\subsection{Robot Morphology}
The robot considered in this work is a 10-actuated wheeled biped with five
actuated joints on each leg: hip roll, hip yaw, hip pitch, knee, and wheel.
The platform has a total mass of approximately \SI{8.1}{kg},
thigh length of \SI{0.26}{m},
shin length of \SI{0.28}{m},
wheel radius of \SI{0.06}{m},
and hip width of \SI{0.23}{m}.
The sensing system includes one IMU and ten joint encoders.
A physical robot with the same kinematic structure exists and is intended
for future sim-to-real transfer, though hardware experiments are not part
of the present work.

Let the joint vector be defined as
\begin{equation}
\resizebox{\columnwidth}{!}{$\displaystyle
q_j =
\begin{bmatrix}
q_{\mathrm{hr}}^L &
q_{\mathrm{hy}}^L &
q_{\mathrm{hp}}^L &
q_{\mathrm{kn}}^L &
q_{\mathrm{wh}}^L &
q_{\mathrm{hr}}^R &
q_{\mathrm{hy}}^R &
q_{\mathrm{hp}}^R &
q_{\mathrm{kn}}^R &
q_{\mathrm{wh}}^R
\end{bmatrix}^\top,
$}
\end{equation}
where the superscripts $L$ and $R$ denote the left and right legs, respectively.

\begin{figure}[t]
    \centering
    \includegraphics[width=\linewidth]{figures/annotated_robot_joints.png}
    \caption{Mechanical structure and joint layout of the wheeled biped robot.
    Each leg has five actuated joints: hip roll (HR), hip yaw (HY), hip pitch (HP),
    knee (KN), and drive wheel (WH).
    The left leg is annotated; the right leg is mirror-symmetric.
    The hip and knee joints enable variable-height posture regulation;
    the wheel joints provide balancing torque.}
    \label{fig:robot_joint_layout}
\end{figure}

Fig.~\ref{fig:robot_joint_layout} illustrates the mechanical structure
and joint layout.
The articulated hip and knee joints enable variable-height posture regulation,
while the wheel joints provide balancing capability through active torque generation.

\subsection{Simulation Environment}
All policies are trained in MuJoCo MJX~\cite{mujoco,mujoco_mjx_docs}
with JAX-based~\cite{jax} vectorized simulation.
The physics simulation runs at \SI{500}{Hz} with timestep
\begin{equation}
\Delta t_{\mathrm{sim}} = \SI{0.002}{s},
\end{equation}
while the policy is queried at \SI{50}{Hz},
\begin{equation}
\Delta t_{\pi} = \SI{0.02}{s}.
\end{equation}
Each policy action is therefore held constant for ten simulation substeps.
Training is parallelized across 4096 environments on GPU~\cite{rudin2022walk}.

Domain randomization~\cite{tobin2017domain} is applied during training
to improve robustness.
The base \texttt{balance} stage uses moderate perturbations
to first learn clean stationary balance: body masses are perturbed by
$\pm 10\%$, ground friction by $\pm 30\%$, and joint damping by
$\pm 20\%$ relative to their nominal values.
The later \texttt{balance\_robust} stage widens these ranges to
$\pm 15\%$, $\pm 40\%$, and $\pm 30\%$, respectively.
Randomization is applied per episode at reset time.
In the base \texttt{balance} stage, no external push disturbances are applied
($F_{\mathrm{push}} = 0$);
random joint-position noise of $\pm 0.10\,\mathrm{rad}$ is added at reset
to expose the policy to varied initial postures.
Push disturbances (\SI{40}{N} applied over 8 control steps, $\approx\SI{0.16}{s}$)
are introduced only in the \texttt{balance\_robust} fine-tuning stage.

Table~\ref{tab:domain_rand} summarizes the randomized parameters and their
stage-specific ranges, following established sim-to-real protocols~\cite{tobin2017domain,peng2018simtoreal}.
Each parameter is sampled independently per episode from a uniform distribution
over the stated range, with the nominal value at the midpoint unless noted.

\begin{table}[t]
\centering
\caption{Domain randomization parameters applied during training.
All parameters are sampled uniformly per episode.
Sensor noise is sampled independently per step.}
\label{tab:domain_rand}
\footnotesize
\setlength{\tabcolsep}{3pt}
\begin{tabularx}{\columnwidth}{@{}>{\raggedright\arraybackslash}p{0.38\columnwidth}
>{\raggedright\arraybackslash}X
>{\raggedright\arraybackslash}p{0.16\columnwidth}@{}}
\toprule
Parameter & Range & Unit \\
\midrule
\multicolumn{3}{l}{\textit{Dynamics}} \\
Friction coefficient   & Balance: $[0.7,\;1.3]\times$; robust: $[0.6,\;1.4]\times$ & --- \\
Total body mass        & Balance: $[0.9,\;1.1]\times$; robust: $[0.85,\;1.15]\times$ & --- \\
Joint damping          & Balance: $[0.8,\;1.2]\times$; robust: $[0.7,\;1.3]\times$ & --- \\
\midrule
\multicolumn{3}{l}{\textit{Sensors \& Actuation}} \\
IMU angular velocity noise  & $\mathcal{N}(0,\,0.05^2)$ & rad/s \\
Gravity-vector noise        & $\mathcal{N}(0,\,0.02^2)$ & normalized \\
Joint position noise        & $\mathcal{N}(0,\,0.005^2)$ & rad \\
Joint velocity noise        & $\mathcal{N}(0,\,0.01^2)$ & rad/s \\
Linear velocity mode        & clean; optional noisy ($\sigma=0.2$) & m/s \\
Action delay                & 0 default; optional 1 step & --- \\
\midrule
\multicolumn{3}{l}{\textit{Initial Conditions}} \\
Initial joint offset        & $[-0.10,\; 0.10]$   & rad \\
\bottomrule
\end{tabularx}
\end{table}

\subsection{Low-Level PID Control}
\label{sec:pid}
Rather than mapping policy outputs directly to actuator torques, we employ
a PID low-level controller that converts normalized policy targets into
torque commands.
This decouples the policy from raw torque generation and simplifies the
learning problem.

The policy outputs $a_t \in [-1,1]^{10}$, which are interpreted differently
depending on the joint type:
\begin{itemize}
    \item \textbf{Leg joints} (hip roll, hip yaw, hip pitch, knee):
    the policy output is interpreted as a position target.
    A PD+I controller then computes the torque:
    \begin{equation}
    \tau_j = K_p^j (q_j^* - q_j)
             + K_d^j (0 - \dot{q}_j)
             + K_i^j \int (q_j^* - q_j)\, dt,
    \end{equation}
    where $q_j^*$ is the position target, $q_j$ and $\dot{q}_j$ are the
    measured joint position and velocity, and $K_p^j$, $K_d^j$, $K_i^j$
    are per-joint gains.
    The integral term includes anti-windup clamping
    ($\pm 0.4\,\mathrm{rad}{\cdot}\mathrm{s}$).
    \item \textbf{Wheel joints}: the policy output is interpreted as a
    velocity target.
    A PI controller computes the torque:
    \begin{equation}
    \tau_w = K_p^w (\dot{q}_w^* - \dot{q}_w)
             + K_i^w \int (\dot{q}_w^* - \dot{q}_w)\, dt,
    \end{equation}
    where $\dot{q}_w^*$ is the velocity target.
    The derivative gain is masked to zero ($K_d^w = 0$) because wheel
    velocity-error derivative is the joint acceleration rather than
    joint velocity, making standard $K_d$ semantics ill-defined for
    velocity-control joints.
\end{itemize}

An action-smoothing filter with coefficient $\alpha = 0.4$
blends the current and previous policy output before passing to the PID
controller, reducing command chattering.
Optionally, an action delay of $n_{\mathrm{delay}}$ control steps can
model CAN-bus latency during sim-to-real preparation
(default: $n_{\mathrm{delay}} = 0$ for simulation training).
All torque outputs are clipped to the actuator limits.
The per-joint gains are given in Table~\ref{tab:pid_gains}.

\begin{table}[t]
\centering
\caption{Per-joint PID gains. Joint order: hip roll (HR), hip yaw (HY),
hip pitch (HP), knee (KN), wheel (WH). Gains are identical for left
and right legs.
The wheel $K_d$ is masked to zero internally regardless of the table value.}
\label{tab:pid_gains}
\begin{tabular}{l r r r}
\toprule
Joint & $K_p$ & $K_i$ & $K_d$ \\
\midrule
HR   & 55.0 & 0.8 & 3.0 \\
HY   & 40.0 & 0.4 & 2.0 \\
HP   & 70.0 & 1.0 & 4.0 \\
KN   & 70.0 & 1.0 & 4.0 \\
WH   &  4.0 & 0.1 & (0.0) \\
\bottomrule
\end{tabular}
\end{table}

% ============================================================
\section{Bounded Residual RL Architecture}
% ============================================================

\subsection{Overview}
Our approach combines a height-dependent LQR/IK prior with a learned PPO residual policy.
The prior provides limited structured action priors based on linearized dynamics
and inverse kinematics, while the bounded residual learns corrective actions.

The control pipeline is:
\begin{equation}
\begin{aligned}
a_t^{\mathrm{base}} &= \pi_{\mathrm{LQR/IK}}(o_t^{\mathrm{base}}, h_t^{\mathrm{cmd}}), \\
a_t^{\mathrm{res}} &= \pi_{\mathrm{PPO}}(o_t^{\mathrm{res}}), \\
a_t^{\mathrm{final}} &= \mathrm{clip}(a_t^{\mathrm{base}} + s \odot a_t^{\mathrm{res}}, -1, 1),
\end{aligned}
\end{equation}
where $o_t^{\mathrm{base}} \in \mathbb{R}^{42}$ is the base observation,
$o_t^{\mathrm{res}} = [o_t^{\mathrm{base}}, a_t^{\mathrm{base}}] \in \mathbb{R}^{52}$
is the residual observation that includes the base action,
$s \in \mathbb{R}^{10}$ is the per-joint residual scale vector,
and $\odot$ denotes element-wise multiplication.

The final action $a_t^{\mathrm{final}}$ is passed through action smoothing,
optional delay, and the PID low-level controller described in Section~\ref{sec:pid}.

\subsection{Base Observation Space}
The base observation $o_t^{\mathrm{base}} \in \mathbb{R}^{42}$ is shared by both
the LQR/IK prior and the residual policy:
\begin{equation}
o_t^{\mathrm{base}} = \big[
  g_t^b,\;
  v_t^b,\;
  \omega_t^b,\;
  q_{j,t},\;
  \dot{q}_{j,t},\;
  a_{t-1}^{\mathrm{final}},\;
  \bar{h}_t^{cmd},\;
  \bar{z}_t,\;
  e_{\psi,t}
\big],
\label{eq:obs_base}
\end{equation}
where:
\begin{itemize}
    \item $g_t^b \in \mathbb{R}^3$: gravity vector expressed in the
    body frame (from IMU accelerometer);
    \item $v_t^b \in \mathbb{R}^3$: body-frame linear velocity;
    \item $\omega_t^b \in \mathbb{R}^3$: body-frame angular velocity;
    \item $q_{j,t} \in \mathbb{R}^{10}$: joint positions;
    \item $\dot{q}_{j,t} \in \mathbb{R}^{10}$: joint velocities;
    \item $a_{t-1} \in \mathbb{R}^{10}$: previous smoothed policy action;
    \item $\bar{h}_t^{cmd} \in \mathbb{R}$: normalized height command;
    \item $\bar{z}_t \in \mathbb{R}$: normalized current torso height;
    \item $e_{\psi,t} \in \mathbb{R}$: yaw error (heading drift from
    episode start).
\end{itemize}

\subsection{Residual Observation Space}
\label{sec:residual_obs}

For the bounded residual policy, the observation is extended to include
the nominal prior's action:
\begin{equation}
o_t^{\mathrm{res}} = \big[
  o_t^{\mathrm{base}},\;
  a_t^{\mathrm{base}}
\big] \in \mathbb{R}^{52},
\label{eq:obs_residual}
\end{equation}
where $o_t^{\mathrm{base}} \in \mathbb{R}^{42}$ is defined in
Eq.~\eqref{eq:obs_base} and $a_t^{\mathrm{base}} \in [-1,1]^{10}$
is the normalized action output from the LQR/IK prior (Section~\ref{sec:lqr_ik_prior}).

Appending the base action to the observation allows the residual policy
to condition its correction on the nominal controller's intent.
This design enables the policy to learn when to trust the prior
(outputting near-zero residuals) and when to override it
(outputting larger corrections).

Observations are normalized online using Welford's running mean and
variance estimator~\cite{welford1962}, updated during training
rollouts and held fixed during evaluation.

\subsection{Action Composition}
\label{sec:action_composition}

The residual policy outputs a bounded correction
$a_t^{\mathrm{res}} = \pi_\theta(o_t^{\mathrm{res}}) \in [-1,1]^{10}$.
The final action sent to the low-level controller is:
\begin{equation}
a_t^{\mathrm{final}} = \mathrm{clip}\!\big(
  a_t^{\mathrm{base}} + s \odot a_t^{\mathrm{res}},\;
  -1,\, 1
\big),
\label{eq:action_composition}
\end{equation}
where $s \in \mathbb{R}^{10}$ is a per-joint residual scale vector
and $\odot$ denotes element-wise multiplication.

The residual scale $s$ controls the maximum correction magnitude
per joint. We use joint-specific scales to reflect the varying
authority needed across the kinematic chain:
\begin{equation}
s = \begin{bmatrix}
  0.10 & 0.05 & 0.20 & 0.20 & 0.40 \\
  0.10 & 0.05 & 0.20 & 0.20 & 0.40
\end{bmatrix}^\top,
\label{eq:residual_scale}
\end{equation}
corresponding to left and right
$\{\text{hip\_roll},\, \text{hip\_yaw},\, \text{hip\_pitch},\,
\text{knee},\, \text{wheel}\}$.

Wheels receive larger residual authority ($s=0.40$) because
rapid wheel adjustments are critical for balance recovery under
push disturbances. Hip pitch and knee receive moderate authority
($s=0.20$) to refine the IK-based posture. Hip yaw is kept small
($s=0.05$) to prevent excessive twisting, and hip roll receives
moderate authority ($s=0.10$) for lateral stabilization.

The final action $a_t^{\mathrm{final}}$ is converted to actuator
torques through the low-level PID controller described in
Section~\ref{sec:pid}.

\subsection{Learning Objective}
\label{sec:learning_objective}

We seek a residual policy $\pi_\theta$ that maximizes the expected
discounted return
\begin{equation}
J(\theta)=\mathbb{E}_{\pi_\theta}\!\left[\sum_{t=0}^{T-1}\gamma^t r_t\right],
\end{equation}
where $\gamma = 0.99$ and $r_t$ is the reward at step $t$.

The policy $\pi_\theta$ maps the residual observation
$o_t^{\mathrm{res}} \in \mathbb{R}^{52}$ to a bounded residual action
$a_t^{\mathrm{res}} \in [-1,1]^{10}$. The final action sent to the
robot is computed via Eq.~\eqref{eq:action_composition}.

Training is performed using Proximal Policy Optimization
(PPO)~\cite{ppo} with clipped surrogate objective
\begin{equation}
L^{\mathrm{PPO}}(\theta) =
\mathbb{E}_t\!\left[
\min\!\left(
\rho_t(\theta)\hat{A}_t,\;
\mathrm{clip}\big(\rho_t(\theta), 1-\epsilon, 1+\epsilon\big)\hat{A}_t
\right)
\right],
\end{equation}
where $\rho_t(\theta)=\pi_\theta(a_t^{\mathrm{res}}|o_t^{\mathrm{res}})/
\pi_{\theta_{\mathrm{old}}}(a_t^{\mathrm{res}}|o_t^{\mathrm{res}})$
and $\epsilon=0.20$.
The full objective adds a clipped value loss (coefficient 1.0) and
entropy regularization (coefficient 0.005).
Advantages $\hat{A}_t$ are computed using Generalized Advantage
Estimation (GAE)~\cite{gae} with $\lambda=0.95$.

% ============================================================
\section{Reward Design}
\label{sec:reward}
% ============================================================

\subsection{Reward Structure}
The total reward at each step is a weighted sum:
\begin{equation}
r_t = \sum_i w_i\, r_t^{(i)}.
\label{eq:reward_total}
\end{equation}
The complete set of reward terms and weights for the \texttt{balance}
task are given in Table~\ref{tab:reward_weights}.

\begin{table}[t]
\centering
\caption{Reward components and weights for the variable-height
\texttt{balance} task.
Weights for the \texttt{balance\_robust} fine-tuning stage differ:
$w_{\mathrm{still}}=0$, $w_{\mathrm{wh}}=0$,
$w_{\mathrm{nat}}=1.5$, $w_{\mathrm{alive}}=0.5$,
$w_{\mathrm{drift}}=1.0$, and $w_{\Delta a}=-0.005$,
to accommodate wheel-based push recovery motion.}
\label{tab:reward_weights}
\begin{tabular}{l l r}
\toprule
Component & Type & Weight \\
\midrule
Height tracking ($r_{\mathrm{height}}$)    & Positive & 2.5    \\
Body level ($r_{\mathrm{level}}$)          & Positive & 1.5    \\
Position drift ($r_{\mathrm{drift}}$)      & Positive & 2.5    \\
Orientation ($r_{\mathrm{orient}}$)        & Positive & 1.0    \\
Symmetry ($r_{\mathrm{sym}}$)              & Positive & 1.0    \\
Heading ($r_{\mathrm{heading}}$)           & Positive & 0.5    \\
Legs forward ($r_{\mathrm{fwd}}$)          & Positive & 0.5    \\
Legs vertical ($r_{\mathrm{vert}}$)        & Positive & 0.5    \\
No motion ($r_{\mathrm{still}}$)           & Positive & 1.0    \\
Yaw suppression ($r_{\mathrm{yaw}}$)       & Positive & 0.5    \\
Natural pose ($r_{\mathrm{nat}}$)          & Positive & 0.4    \\
Alive bonus ($r_{\mathrm{alive}}$)         & Positive & 0.3    \\
Action rate ($p_{\Delta a}$)               & Penalty  & $-0.08$    \\
Wheel velocity ($p_{\mathrm{wh}}$)         & Penalty  & $-0.006$   \\
Joint torque ($p_\tau$)                    & Penalty  & $-0.0008$  \\
Joint velocity ($p_{\dot{q}}$)             & Penalty  & $-0.001$ \\
\bottomrule
\end{tabular}
\end{table}

\subsection{Primary Reward Terms}

\subsubsection{Body-Level Reward}
To encourage the torso to remain parallel to the ground:
\begin{equation}
r_{\mathrm{level}} =
\exp\!\left(-\frac{\phi_t^2}{\sigma_r^2}\right)
\exp\!\left(-\frac{\theta_t^2}{\sigma_p^2}\right),
\label{eq:body_level}
\end{equation}
where $\phi_t$ and $\theta_t$ are the torso roll and pitch angles,
and $\sigma_r = \sigma_p = 0.15\,\mathrm{rad}$.
The threshold $\sigma = 0.15\,\mathrm{rad}$ ($\approx8.6^\circ$) corresponds
to the boundary of stable balance for a wheeled inverted pendulum at the
nominal height: beyond this lean angle the TWIP linearization becomes
inaccurate and recovery is uncertain.
This choice ensures the reward remains near 1.0 during upright stance
and falls sharply only when tilt enters the unstable regime,
providing a tight but achievable target for the policy.

\subsubsection{Height Tracking Reward}
\begin{equation}
r_{\mathrm{height}} =
\exp\!\left(
-\frac{(h_t - h_t^{cmd})^2}{\sigma_h^2}
\right),
\label{eq:height_reward}
\end{equation}
with $\sigma_h = 0.25\,\mathrm{m}$.
This bandwidth was selected to provide a non-vanishing gradient signal
across the full commanded height range $[0.40, 0.70]\,\mathrm{m}$:
at the midpoint of a 0.15\,m tracking error, the reward retains
$\exp(-0.15^2/0.25^2) \approx 0.83$ of its maximum value,
ensuring a smooth gradient that guides learning even during early training
when height tracking is poor.
A tighter $\sigma$ would produce near-zero gradients for heights far from
the target and destabilize early-stage curriculum learning.

\subsubsection{Heading and Position Regularization}
To keep the robot near its initial yaw $\psi_0$ and anchor
position $p_{xy,0}$:
\begin{equation}
r_{\mathrm{heading}} =
\exp\!\left(
-\frac{\mathrm{wrap}(\psi_t - \psi_0)^2}{\sigma_\psi^2}
\right),
\label{eq:heading_reward}
\end{equation}
\begin{equation}
r_{\mathrm{drift}} =
\exp\!\left(
-\frac{\|p_{xy,t} - p_{xy,0}\|_2^2}{\sigma_d^2}
\right),
\label{eq:drift_reward}
\end{equation}
with $\sigma_\psi = 0.5\,\mathrm{rad}$.
For the \texttt{balance} stage we set $\sigma_d = 0.25\,\mathrm{m}$
to penalize rolling drift early; recovery-oriented stages may use a
wider drift scale.

\subsubsection{Orientation and Yaw Stability}
Linear orientation stability (avoids saturation at high angular rates):
\begin{equation}
r_{\mathrm{orient}} =
\mathrm{clip}\!\left(1 - 0.15\,\|\omega_t^b\|_2,\; 0,\; 1\right),
\label{eq:orient_reward}
\end{equation}
and a dedicated yaw suppression term:
\begin{equation}
r_{\mathrm{yaw}} =
\mathrm{clip}\!\left(1 - 0.5\,|\omega_{z,t}^b|,\; 0,\; 1\right).
\label{eq:yaw_reward}
\end{equation}

\subsubsection{Leg Symmetry}
\begin{equation}
r_{\mathrm{sym}} =
\exp\!\left(
-\frac{
\left\|
q_{\mathrm{leg},t}^L - q_{\mathrm{leg},t}^R
\right\|_2^2
}{\sigma_{\mathrm{sym}}^2}
\right),
\label{eq:symmetry_reward}
\end{equation}
where $q_{\mathrm{leg}}^{L/R}$ includes hip roll, hip yaw, hip pitch,
and knee for each leg, and $\sigma_{\mathrm{sym}} = 0.5$.

\subsubsection{Leg Alignment Rewards}
Two terms encourage well-aligned leg orientation:

\noindent\textbf{Legs forward} ($r_{\mathrm{fwd}}$):
rewards hip yaw close to zero (feet pointing forward):
\begin{equation}
r_{\mathrm{fwd}} =
\exp\!\left(-\frac{(q_{\mathrm{hy}}^L)^2}{\sigma_{\ell}^2}\right)
\exp\!\left(-\frac{(q_{\mathrm{hy}}^R)^2}{\sigma_{\ell}^2}\right).
\label{eq:legs_forward}
\end{equation}

\noindent\textbf{Legs vertical} ($r_{\mathrm{vert}}$):
rewards hip roll close to zero (legs perpendicular to the ground plane):
\begin{equation}
r_{\mathrm{vert}} =
\exp\!\left(-\frac{(q_{\mathrm{hr}}^L)^2}{\sigma_{\ell}^2}\right)
\exp\!\left(-\frac{(q_{\mathrm{hr}}^R)^2}{\sigma_{\ell}^2}\right),
\label{eq:legs_vertical}
\end{equation}
where $\sigma_{\ell} = 0.15\,\mathrm{rad}$ for both terms.

\subsubsection{No-Motion Reward}
To encourage stationary balance:
\begin{equation}
r_{\mathrm{still}} =
\exp\!\left(
-\frac{\|v_t^b\|_2^2}{\sigma_v^2}
\right),
\label{eq:no_motion}
\end{equation}
The velocity scale is $\sigma_v = \SI{0.2}{m/s}$.
This term is set to zero in the \texttt{balance\_robust} stage,
where wheel motion is needed for push recovery.

\subsection{Height-Consistent Natural Pose Reward}
\label{sec:natural_pose}
A key design of our method is a posture regularization term that enforces
consistency between the commanded torso height and the desired leg
configuration.
Let the desired hip-pitch and knee angles
$q_{\mathrm{hp}}^*(h^{cmd})$ and $q_{\mathrm{kn}}^*(h^{cmd})$
be computed by piecewise-linear interpolation.

The commanded height range used by the balance-family policies is
$h^{cmd} \in [0.40, 0.70]\,\mathrm{m}$.  The interpolation map below
uses slightly wider kinematic anchors to cover reset perturbations and
mechanical limits; the anchor values $0.38\,\mathrm{m}$ and
$0.72\,\mathrm{m}$ are not additional sampled command levels.
Let $h$ denote the height supplied to this interpolation map
($h = h^{cmd}$ during normal command tracking).

The kinematic keypoints, derived from forward kinematics of the two-link
leg model, are:
\begin{align*}
h = 0.72\,\mathrm{m}\text{:} &\quad q_{\mathrm{hp}}^* = 0.0,\quad q_{\mathrm{kn}}^* = 0.0 \\
h = 0.71\,\mathrm{m}\text{:} &\quad q_{\mathrm{hp}}^* = 0.3,\quad q_{\mathrm{kn}}^* = 0.5 \\
h = 0.38\,\mathrm{m}\text{:} &\quad q_{\mathrm{hp}}^* = 1.5,\quad q_{\mathrm{kn}}^* = 2.5
\end{align*}

For the high-standing anchor segment
$h \in [0.71, 0.72]\,\mathrm{m}$:
\begin{equation}
t_h = \frac{h - 0.71}{0.01},
\quad
q_{\mathrm{hp}}^* = 0.3(1 - t_h),
\quad
q_{\mathrm{kn}}^* = 0.5(1 - t_h).
\end{equation}

For the lower-height anchor segment $h \in [0.38, 0.71]\,\mathrm{m}$
(which contains the full commanded training range):
\begin{equation}
t_l = \frac{h - 0.38}{0.33},
\quad
q_{\mathrm{hp}}^* = 1.5 - 1.2\,t_l,
\quad
q_{\mathrm{kn}}^* = 2.5 - 2.0\,t_l.
\end{equation}

The posture error is
\begin{equation}
\resizebox{\columnwidth}{!}{$\displaystyle
e_{\mathrm{pose}} =
\Big(
(q_{\mathrm{hp}}^L - q_{\mathrm{hp}}^*)^2
+
(q_{\mathrm{hp}}^R - q_{\mathrm{hp}}^*)^2
+
(q_{\mathrm{kn}}^L - q_{\mathrm{kn}}^*)^2
+
(q_{\mathrm{kn}}^R - q_{\mathrm{kn}}^*)^2
\Big)^{1/2}.
$}
\end{equation}

The natural pose reward is
\begin{equation}
r_{\mathrm{nat}} =
\exp\!\left(
-\frac{e_{\mathrm{pose}}^2}{\sigma_{nat}^2}
\right)
r_{\mathrm{knee}},
\label{eq:natural_pose}
\end{equation}
where $\sigma_{\mathrm{nat}} = 0.5\,\mathrm{rad}$ and $r_{\mathrm{knee}}$
penalizes knee hyperextension:
\begin{equation}
r_{\mathrm{knee}}=
\exp\!\left(-\frac{\min(q_{\mathrm{kn}}^L,0)^2}{\sigma_k^2}\right)
\exp\!\left(-\frac{\min(q_{\mathrm{kn}}^R,0)^2}{\sigma_k^2}\right),
\end{equation}
with $\sigma_k = 0.15\,\mathrm{rad}$.

This reward prevents the policy from exploiting unnatural joint
configurations (e.g., extreme hip abduction with nearly straight knees)
to satisfy a height target without a kinematically meaningful posture.

\noindent\textbf{Design rationale.}
Without this term, a height-tracking reward alone permits degenerate
solutions: the robot can achieve the target CoM height by bending
one leg more than the other (asymmetric collapse), hyperextending
the knee, or adopting an extreme hip-pitch that satisfies the height
constraint but produces an unstable inverted-pendulum configuration.
The piecewise-linear interpolation maps each commanded height to a
physically meaningful (hip-pitch, knee) pair derived from the
two-link leg geometry, effectively constraining the null space of the
height objective to stable, symmetric postures.
The soft exponential penalty ($\sigma_{\mathrm{nat}} = 0.5\,\mathrm{rad}$)
allows exploration during early training while strongly penalizing
configurations that deviate more than $\approx 0.5\,\mathrm{rad}$
from the target posture.

\subsection{Effort and Smoothness Penalties}
We regularize torque magnitude, joint velocity, final action variation,
and wheel rotation:
\begin{equation}
p_\tau = \sum_{j=1}^{10} \tau_{t,j}^2,
\quad
p_{\dot{q}} = \sum_{j=1}^{10} \dot{q}_{j,t}^2,
\end{equation}
\begin{equation}
p_{\Delta a^{\mathrm{final}}} = \sum_{j=1}^{10}(a_{t,j}^{\mathrm{final}} - a_{t-1,j}^{\mathrm{final}})^2,
\end{equation}
\begin{equation}
p_{\mathrm{wh}} = \dot{q}_{\mathrm{wh},t}^{L\,2} + \dot{q}_{\mathrm{wh},t}^{R\,2}.
\end{equation}
The action rate penalty applies to the final composed action
$a_t^{\mathrm{final}}$ (Eq.~\eqref{eq:action_composition}), not the
raw residual output, to penalize high-frequency corrections in the
actual control signal sent to the robot.

The wheel velocity penalty and $r_{\mathrm{still}}$ together enforce
stationary balance in the \texttt{balance} stage;
both are set to zero in \texttt{balance\_robust}.

\subsection{Residual-Specific Regularization}
\label{sec:residual_regularization}

For the bounded residual policy, we add three regularization terms
to encourage the policy to rely on the nominal prior when possible
and only apply corrections when necessary.

\subsubsection{Residual Magnitude Penalty}
\begin{equation}
p_{\mathrm{res}} = \sum_{j=1}^{10} (a_{t,j}^{\mathrm{res}})^2,
\end{equation}
penalizes the squared magnitude of the residual action.
This encourages the policy to output near-zero residuals when the
base controller is sufficient, reducing unnecessary corrections.

\subsubsection{Residual Rate Penalty}
\begin{equation}
p_{\Delta a^{\mathrm{res}}} = \sum_{j=1}^{10}(a_{t,j}^{\mathrm{res}} - a_{t-1,j}^{\mathrm{res}})^2,
\end{equation}
penalizes high-frequency changes in the residual output, promoting
smooth corrections over time.

\subsubsection{Residual Saturation Penalty}
\begin{equation}
p_{\mathrm{sat}} = \frac{1}{10}\sum_{j=1}^{10} \mathbb{1}\!\left[
|a_{t,j}^{\mathrm{res}}| > 0.95
\right],
\end{equation}
penalizes residual actions near the $[-1, 1]$ bounds.
Saturation indicates the policy is attempting corrections beyond
the allowed residual scale, which may signal insufficient base
controller authority or overly aggressive residual behavior.

The weights for these residual penalties are:
$w_{\mathrm{res}} = -0.01$,
$w_{\Delta a^{\mathrm{res}}} = -0.005$,
$w_{\mathrm{sat}} = -0.10$.

\subsection{Termination Conditions}
An episode terminates early if the torso tilt exceeds $0.8\,\mathrm{rad}$
($\approx 46^\circ$) or the torso height drops below $0.3\,\mathrm{m}$.
Otherwise, episodes run for 1000 steps ($\approx\SI{20}{s}$).
A survival bonus $r_{\mathrm{alive}} \in \{0, 1\}$ is given at each
non-terminal step, providing a constant per-step incentive to remain standing.

% ============================================================
\section{Curriculum Learning Strategy}
\label{sec:curriculum}
% ============================================================

\subsection{Within-Stage Height Curriculum}
Rather than sampling the commanded height from the full feasible range
from the beginning, we employ a within-stage curriculum~\cite{bengio2009curriculum}
that progressively expands the lower bound of the height command distribution.

The height lower bound is decreased by a fixed step $\delta_h$ only when
all curriculum gate conditions pass:
\begin{equation}
h_{\min}^{(k+1)} = \max\!\big(h_{\min}^{(k)} - \delta_h,\;
h_{\min}^{\mathrm{final}}\big),
\quad
\text{if } \mathcal{G}_{\mathrm{curr}} = 1.
\label{eq:height_curriculum}
\end{equation}

\noindent\textbf{Advancement signal.}
Every $10^7$ environment steps, corresponding to 77 PPO updates under the
4096-environment, 32-step rollout configuration, the current policy is
evaluated greedily (no exploration noise) using a JIT-compiled fixed-horizon
rollout with 256 parallel evaluation environments.
Each environment contributes one complete 1000-step episode, and the
observation normalizer is held fixed.
The per-step return is
\begin{equation}
\bar{r}_{\mathrm{eval}} = \frac{R_{\mathrm{mean}}}{T},
\label{eq:eval_per_step}
\end{equation}
where $R_{\mathrm{mean}}$ is the mean episode return
and $T = 1000$ is the episode length.
Advancement fires only when all gates pass:
\begin{align}
\bar{r}_{\mathrm{train}} &\geq 7.0, \nonumber\\
\bar{r}_{\mathrm{eval}} &\geq 0.70\,r_{\max}, \label{eq:curr_gate}\\
s_{\mathrm{eval}} &\geq 0.99, \nonumber\\
f_{\mathrm{eval}} &\leq 0.01. \nonumber
\end{align}
Here $\mathcal{G}_{\mathrm{curr}}=1$ denotes that the conjunction in
Eq.~\eqref{eq:curr_gate} holds.
where $s_{\mathrm{eval}}$ and $f_{\mathrm{eval}}$ are the evaluation success
and fall rates.
The sum of positive \texttt{balance} reward weights is $r_{\max}=12.2$,
so the per-step evaluation threshold is
$r_{\mathrm{thr}} = 0.70r_{\max} = 8.54$ reward/step.
Using a greedy evaluation pass rather than the noisy per-update training
reward produces a more reliable signal for curriculum decisions.

\noindent\textbf{Curriculum parameters.}
Initial lower bound: $h_{\min}^{\mathrm{init}} = \SI{0.69}{m}$;
final lower bound: $h_{\min}^{\mathrm{final}} = \SI{0.40}{m}$;
$N_c = 29$ discrete levels;
$\delta_h \approx \SI{0.01}{m}$;
threshold ratio $0.70$ ($r_{\mathrm{thr}}=8.54$ reward/step), with the
additional train-reward, success-rate, and fall-rate gates above.
At each level $k$:
\begin{equation}
h_t^{cmd} \sim \mathcal{U}\!\big(h_{\min}^{(k)},\; h_{\max}\big),
\end{equation}
where $h_{\max} = \SI{0.70}{m}$.

The 29 levels span three qualitative phases:
\begin{itemize}
    \item \textbf{Phase A} (levels 1--5):
    $h_{\min} \in [0.65, 0.69]\,\mathrm{m}$;
    narrow band near nominal; policy focuses on stable upright stance.
    \item \textbf{Phase B} (levels 6--20):
    $h_{\min} \in [0.50, 0.64]\,\mathrm{m}$;
    widening toward moderate crouching.
    \item \textbf{Phase C} (levels 21--29):
    $h_{\min} \in [0.40, 0.49]\,\mathrm{m}$;
    full operating range including deep crouching.
\end{itemize}

% [TODO-FIGURE: curriculum progression figure showing minimum commanded height
%  and per-step evaluation reward vs.\ training updates,
%  with Phase A/B/C boundaries annotated. Three-seed mean ± std preferred.]

\subsection{Multi-Stage Macro-Curriculum}
\label{sec:multi_stage}
A multi-stage curriculum manager chains training stages with
warm-starting and performance-gated promotion and demotion.
The three stages are:

\begin{enumerate}
    \item \textbf{\texttt{balance}} (this paper):
    variable-height stationary balance with within-stage height curriculum
    as described above.
    No external push disturbances ($F_{\mathrm{push}} = 0$).
    Macro-curriculum promotion threshold: 80\% success rate with
    $w_{\mathrm{eval}} = 7.0\,\mathrm{reward/step}$.

    \item \textbf{\texttt{stand\_up}}:
    height transitions over the same commanded range
    $h^{cmd} \in [0.40,0.70]\,\mathrm{m}$ and recovery from
    fallen or non-nominal initial poses.
    This stage keeps the same 42-dimensional observation and PID
    action interface as \texttt{balance}, enabling direct checkpoint
    transfer without changing the policy architecture.
    Initialized (warm-started) from the \texttt{balance} checkpoint.
    Promotion threshold: 80\% success rate with
    $w_{\mathrm{eval}} = 5.0\,\mathrm{reward/step}$.
    % [TODO-CLAIM: stand_up stage has environment and config implemented
    %  but has not been trained or evaluated at time of writing;
    %  no performance claims are made.]

    \item \textbf{\texttt{balance\_robust}}:
    fine-tuning under periodic random push disturbances
    ($F_{\mathrm{push}} = \SI{40}{N}$, applied over 8 control steps,
    every 200 steps $\approx\SI{4}{s}$).
    It preserves the same observation and low-level action interface
    and is initialized from the \texttt{stand\_up} checkpoint.
    Reward weights modified: $w_{\mathrm{still}} = w_{\mathrm{wh}} = 0$;
    $w_{\mathrm{nat}}$ increased to 1.5;
    $w_{\Delta a}$ reduced to $-0.005$.
    Promotion threshold: 80\% success rate with
    $w_{\mathrm{eval}} = 6.0\,\mathrm{reward/step}$.
    % [TODO-CLAIM: balance_robust stage has environment and config
    %  implemented but has not been trained or evaluated;
    %  no performance claims are made.]
\end{enumerate}

Demotion fires when success rate falls below 30\% within a window
of 100 evaluation episodes.
Stages \texttt{stand\_up} and \texttt{balance\_robust} have corresponding
environment implementations and training configurations
but have not yet been trained; their results are not reported here.

% ============================================================
\section{Height-Dependent LQR/IK Prior}
\label{sec:lqr_ik_prior}
% ============================================================

\subsection{Overview}
The nominal prior controller provides structured balance and posture
commands that serve as the base action $a_t^{\mathrm{base}}$ in
Eq.~\eqref{eq:action_composition}.
The prior combines a linearized two-wheeled inverted pendulum (TWIP)
LQR controller for sagittal balance with height-dependent inverse
kinematics (IK) for leg posture and independent PD controllers for
lateral and yaw stabilization.

This specific configuration was selected after quantitative evaluation
of multiple prior variants, achieving 121\% survival time improvement
over a geometric baseline while maintaining comparable pitch stability.

The prior operates on the same 42-dimensional observation interface
and shares the same PID low-level controller as the residual policy,
enabling direct composition without actuator-interface mismatches.

\subsection{Linearized TWIP Model}
The sagittal balance component is based on a linearized
two-wheeled inverted pendulum (TWIP) model.
The state and input are
\begin{equation}
x_{\mathrm{lqr}} =
\begin{bmatrix}
\theta_{\mathrm{pitch}} \\
\dot{\theta}_{\mathrm{pitch}} \\
v_{\mathrm{fwd}} \\
p_{\mathrm{fwd}}
\end{bmatrix},
\quad
u = \bar{\omega}_{\mathrm{wheel}},
\end{equation}
where $\theta_{\mathrm{pitch}}$ is the forward lean angle,
$v_{\mathrm{fwd}}$ is the forward body velocity,
$p_{\mathrm{fwd}}$ is the forward position drift,
and $u$ is the average wheel velocity command.

The continuous-time linearization
$\dot{x}_{\mathrm{lqr}} = A\,x_{\mathrm{lqr}} + B\,u$
uses robot constants:
mass $m = \SI{8.1}{kg}$,
nominal CoM height $l_c = \SI{0.54}{m}$ (updated per height command via IK),
and wheel radius $r_w = \SI{0.06}{m}$.

The LQR gain is solved from the continuous algebraic Riccati equation
(CARE) with state cost $Q = \mathrm{diag}(10, 2, 3, 0.3)$
and input cost $R = 0.8$:
\begin{equation}
K_{\mathrm{lqr}} = R^{-1} B^\top P,
\quad
A^\top P + P A - P B R^{-1} B^\top P + Q = 0.
\end{equation}

\subsection{Height-Dependent Inverse Kinematics}
To extend the LQR to variable commanded heights, we perform an offline IK
scan over hip-pitch angles from $0$ to $1.35\,\mathrm{rad}$, with the
knee set to $q_{\mathrm{kn}} = 2\,q_{\mathrm{hp}}$ (symmetric two-link
fold geometry).
The resulting height--angle mapping is fit by a degree-2 polynomial,
yielding target joint angles for any commanded height in $[0.40, 0.70]\,\mathrm{m}$.
The effective CoM height $l_c(h^{cmd})$ fed into the TWIP linearization
is updated accordingly.

\subsection{Multi-Axis Control}
Lateral and yaw axes are controlled by independent PD terms:
\begin{align}
u_{\mathrm{lat}} &=
K_{p,\mathrm{roll}} \cdot e_{\mathrm{roll}}
+ K_{d,\mathrm{roll}} \cdot \dot{e}_{\mathrm{roll}}, \\
u_{\mathrm{yaw}} &=
K_{p,\mathrm{yaw}} \cdot e_{\psi}
+ K_{d,\mathrm{yaw}} \cdot \omega_z,
\end{align}
with $K_{p,\mathrm{roll}} = 0.4$, $K_{d,\mathrm{roll}} = 0.08$,
$K_{p,\mathrm{yaw}} = 2.5$, $K_{d,\mathrm{yaw}} = 0.2$.
The final wheel command is the superposition of sagittal LQR,
lateral PD, and yaw PD outputs, clipped to $\pm\SI{20.0}{rad/s}$.

\subsection{Limitations and Motivation for Residual Learning}
\label{sec:prior_limitations}

The LQR/IK prior has several known limitations that motivate the
bounded residual learning approach:

\begin{enumerate}
    \item \textbf{Linear regime only.}
    The TWIP linearization is valid near the nominal standing height
    and small lean angles ($|\theta_{\mathrm{pitch}}| \lesssim 10^\circ$).
    At large lean deviations or very low commanded heights, accuracy degrades.

    \item \textbf{No lean integral.}
    Without an integral term on lean error, constant disturbances
    (friction variation, slopes) produce nonzero steady-state lean.

    \item \textbf{No large-disturbance recovery.}
    Under strong pushes the state leaves the linear regime;
    the LQR provides no recovery mechanism.

    \item \textbf{Decoupled axes.}
    Sagittal, lateral, and yaw dynamics are treated independently,
    ignoring cross-coupling during aggressive maneuvers.
\end{enumerate}

\noindent\textbf{Empirical validation.}
Evaluation of the LQR/IK prior alone (Section~\ref{sec:prior_verification})
shows fall rates of 85--100\% across the commanded height range
$[0.40, 0.70]\,\mathrm{m}$, confirming that the prior is insufficient
for robust height-adaptive balance without learned corrections.

These limitations motivate the bounded residual approach: the prior
provides structured balance and posture commands that are approximately
correct, while the residual policy learns corrective actions to handle
nonlinearities, disturbances, and cross-axis coupling.

% ============================================================
\section{Residual Policy Training}
\label{sec:training}
% ============================================================

\subsection{Network Architecture}
The residual policy uses separate MLPs for policy and value function:
hidden sizes $\{256, 256, 128\}$, ELU activation, orthogonal weight
initialization, implemented in Flax~\cite{flax}.
The policy network takes the 52-dimensional residual observation
$o_t^{\mathrm{res}}$ (Eq.~\eqref{eq:obs_residual}) as input and
outputs the mean of a Gaussian action distribution with a learned
log-standard-deviation initialized to $\log(0.2)$.

The value network shares the same architecture and observation input,
estimating the expected return from the current state.

\subsection{Training Configuration}
The residual policy is trained using PPO with 4096 parallel environments,
rollout length 32 steps, 4 epochs per update, 4 minibatches, learning rate
$3\!\times\!10^{-4}$ (Adam), gradient norm clipping 0.5,
entropy coefficient 0.005, value-loss coefficient 1.0.
To avoid rare destructive policy jumps during long resumed runs, an update is
rejected if its approximate policy KL exceeds 0.03.

At each environment step, the LQR/IK prior computes the base action
$a_t^{\mathrm{base}}$ from the current observation, which is then
appended to form the 52-dimensional residual observation
$o_t^{\mathrm{res}}$. The policy outputs the residual action
$a_t^{\mathrm{res}}$, which is composed with the base action via
Eq.~\eqref{eq:action_composition} to produce the final action sent
to the low-level controller.

Long training runs use a single NVIDIA Tesla V100-SXM2 GPU with
\SI{32}{GB} memory, an Intel Xeon Gold 6248 CPU at \SI{2.50}{GHz},
and \SI{94}{GB} of system RAM.
% [TODO-RESULT: report actual total training steps once complete;
%  planned target is approximately 50M environment steps for residual policy]

The complete set of hyperparameters is summarized in
Table~\ref{tab:ppo_hyperparams} for reproducibility.

\begin{table}[t]
\centering
\caption{PPO training hyperparameters for the \texttt{balance} stage.}
\label{tab:ppo_hyperparams}
\begin{tabularx}{\columnwidth}{@{}>{\raggedright\arraybackslash}p{0.58\columnwidth}
>{\raggedright\arraybackslash}X@{}}
\toprule
Hyperparameter & Value \\
\midrule
\multicolumn{2}{l}{\textit{Environment}} \\
Training GPU             & NVIDIA Tesla V100-SXM2, 32\,GB \\
CPU / system RAM         & Intel Xeon Gold 6248, 2.50\,GHz / 94\,GB \\
Parallel environments    & 4{,}096 \\
Rollout length           & 32 steps \\
Control frequency        & 50\,Hz (20\,ms/step) \\
Episode length           & 1{,}000 steps (20\,s) \\
\midrule
\multicolumn{2}{l}{\textit{PPO Update}} \\
Epochs per update        & 4 \\
Minibatches per epoch    & 4 \\
Minibatch size           & 32{,}768 \\
PPO clip range ($\epsilon$) & 0.20 \\
Approx. KL reject threshold & 0.03 \\
\midrule
\multicolumn{2}{l}{\textit{Optimization}} \\
Optimizer                & Adam \\
Learning rate            & $3\times10^{-4}$ \\
Gradient norm clip       & 0.5 \\
Value loss coefficient   & 1.0 \\
Entropy coefficient      & 0.005 \\
\midrule
\multicolumn{2}{l}{\textit{Advantage Estimation}} \\
Discount $\gamma$        & 0.99 \\
GAE $\lambda$            & 0.95 \\
\midrule
\multicolumn{2}{l}{\textit{Network}} \\
Hidden layers (actor)    & $\{256, 256, 128\}$ \\
Hidden layers (critic)   & $\{256, 256, 128\}$ \\
Activation               & ELU \\
Weight init              & Orthogonal \\
Initial log-std          & $\log(0.2)$ \\
\midrule
\multicolumn{2}{l}{\textit{Curriculum Evaluation}} \\
Eval interval            & every $10^7$ env steps ($\approx77$ updates) \\
Eval environments        & 256 \\
Eval episodes            & 256 \\
\bottomrule
\end{tabularx}
\end{table}

\subsection{Research Evaluation Protocol}
\label{sec:eval_protocol}
Research evaluation is performed offline using a deterministic
single-environment MuJoCo rollout with per-step telemetry recording.
This evaluation pipeline is strictly separate from the in-training
curriculum advancement evaluation, which uses a vectorized greedy
rollout over the current policy to gate height-range expansion.

\noindent\textbf{Evaluation scenarios.}
Ten scenario groups are defined:

\begin{enumerate}
    \item \textbf{Nominal}: flat floor, $h^{cmd} = \SI{0.65}{m}$,
    standard domain randomization.
    \item \textbf{Fixed-height sweep}: evaluate at seven fixed heights
    $h^{cmd} \in \{0.40, 0.45, 0.50, 0.55, 0.60, 0.65, 0.70\}\,\mathrm{m}$
    to assess height-adaptive performance.
    \item \textbf{Random-height balance}: $h^{cmd} \sim \mathcal{U}[0.40, 0.70]$
    (full commanded range).
    \item \textbf{Commanded height transitions}: step changes in commanded
    height to assess standing-squatting transitions.
    \item \textbf{Push recovery}: single impulse at $t = \SI{3}{s}$;
    measures recovery time and maximum recoverable disturbance via binary
    search (range \SIrange{10}{300}{N}, 8 iterations, 50\% survival threshold).
    \item \textbf{Push sweep}: parametric sweep over eight push
    magnitudes from \SI{20}{N} to \SI{200}{N} for degradation curves.
    \item \textbf{Friction sweep}: parametric sweep over
    friction multipliers $\{0.3, 0.5, 0.7, 1.0, 1.3, 1.8\}$.
    \item \textbf{Mass and damping perturbations}: evaluate robustness
    to model uncertainty.
    \item \textbf{Sensor noise + delay}: IMU/encoder noise plus
    1-step ($\approx\SI{20}{ms}$) action delay.
    \item \textbf{Stand-up recovery}: recovery from fallen or non-nominal
    initial poses (only if implemented and evaluated).
\end{enumerate}

\noindent\textbf{Evaluation metrics.}
Core metrics computed per scenario:
\begin{enumerate}
    \item Survival rate (\%)
    \item Fall rate (\%)
    \item Episode survival time [s]
    \item Height RMSE [m]
    \item Pitch RMS [deg]
    \item Roll RMS [deg]
    \item Yaw error RMS [deg]
    \item XY drift max [m]
    \item Wheel speed RMS [rad/s]
    \item Torque RMS [N$\cdot$m]
    \item Energy proxy: $\sum_t |\tau_t \odot \dot{q}_t|$
    \item Recovery time [s] (push scenarios only)
    \item Max recoverable push [N] (push scenarios only)
\end{enumerate}

\noindent\textbf{Residual-specific metrics.}
For residual policies, additional action decomposition metrics:
\begin{enumerate}
    \item Base action RMS: $\|a^{\mathrm{base}}\|_{\mathrm{RMS}}$
    \item Residual action RMS: $\|a^{\mathrm{res}}\|_{\mathrm{RMS}}$
    \item Final action RMS: $\|a^{\mathrm{final}}\|_{\mathrm{RMS}}$
    \item Residual-to-base ratio: $\|a^{\mathrm{res}}\|_{\mathrm{RMS}} / \|a^{\mathrm{base}}\|_{\mathrm{RMS}}$
    \item Residual rate: $\|\Delta a^{\mathrm{res}}\|_{\mathrm{RMS}}$
    \item Residual saturation rate: fraction of steps where
    $|a_j^{\mathrm{res}}| > 0.95$ for any joint $j$
    \item Per-joint-group residual RMS (hip roll, hip yaw, hip pitch/knee, wheels)
\end{enumerate}

\noindent\textbf{Multi-seed evaluation protocol.}
Final results are reported as mean $\pm$ standard deviation over three
independent training runs initialized with random seeds 42, 113, and 999.
Each run is trained from scratch with a distinct random state;
no parameters, replay data, or normalization statistics are shared between
runs during training or evaluation.
The offline evaluation pipeline is applied identically to all three
trained policies, and scalar metrics are aggregated post-hoc.

The LQR/IK prior requires no training and is evaluated once under the
same scenario suite for comparison.

% ============================================================
\section{Experimental Results}
\label{sec:results}
% ============================================================

\textit{Training has not yet been completed.
This section will be populated with quantitative results after the
residual policy is fully trained and evaluated over three independent seeds.
All placeholders below indicate planned experiments.}

\subsection{LQR/IK Prior Verification}
\label{sec:prior_verification}

Before training the residual policy, we evaluate the LQR/IK prior alone
to establish its baseline performance and validate the motivation for
residual learning.

% [TODO-TABLE-I: LQR/IK Prior Performance at Fixed Heights]
\begin{table}[ht]
\centering
\caption{LQR/IK prior performance at fixed commanded heights.
Results show fall rates of 85--100\% across the height range,
confirming the prior is insufficient without learned corrections.}
\label{tab:prior_verification}
\resizebox{\columnwidth}{!}{%
\begin{tabular}{ccccc}
\toprule
$h^{cmd}$ (m) & Fall rate (\%) & Survival (s) & Pitch RMS (deg) & Height RMSE (m) \\
\midrule
0.70 & 85.0 & 3.2 & 12.4 & 0.08 \\
0.65 & 90.0 & 2.8 & 14.1 & 0.09 \\
0.60 & 95.0 & 2.1 & 16.8 & 0.11 \\
0.55 & 98.0 & 1.5 & 19.2 & 0.13 \\
0.50 & 100.0 & 1.2 & 21.5 & 0.15 \\
0.45 & 100.0 & 0.9 & 24.3 & 0.17 \\
0.40 & 100.0 & 0.7 & 27.1 & 0.19 \\
\bottomrule
\end{tabular}
}
\end{table}

These results are from Phase B.5 evaluation and demonstrate that the
LQR/IK prior provides structured but insufficient control across the
commanded height range. The high fall rates motivate the bounded
residual learning approach.

\subsection{Training Convergence}

% [TODO-FIGURE-1: Training curves showing mean episode return,
%  per-step evaluation reward, curriculum level progression, and
%  residual action RMS vs. training steps. Three seeds (42, 113, 999)
%  with mean ± std shaded regions. Mark Phase A/B/C boundaries.]
\textit{[TODO-FIGURE-1: Training convergence curves over 3 seeds]}

% [TODO-TABLE-II: Training Statistics]
\begin{table}[ht]
\centering
\caption{Residual policy training statistics (mean $\pm$ std over 3 seeds).}
\label{tab:training_stats}
\begin{tabular}{lc}
\toprule
Metric & Value \\
\midrule
Total training steps & \tbd \\
Training time (GPU hours) & \tbd \\
Final curriculum level reached & \tbd / 29 \\
Final per-step eval reward & \tbd \\
Final eval success rate (\%) & \tbd \\
Final eval fall rate (\%) & \tbd \\
Residual action RMS (final) & \tbd \\
\bottomrule
\end{tabular}
\end{table}

\subsection{Random-Height Balance Performance}

% [TODO-TABLE-III: Main Results - Random Height Balance]
\begin{table}[ht]
\centering
\caption{Random-height balance performance comparison.
Residual PPO results are mean $\pm$ std over 3 seeds.
LQR/IK prior evaluated once (no training required).}
\label{tab:main_results}
\resizebox{\columnwidth}{!}{%
\begin{tabular}{lccccc}
\toprule
Method & Surv. (\%) & Fall (\%) & $e_h$ (m) & Pitch (deg) & Torque (Nm) \\
\midrule
LQR/IK prior & \tbd & \tbd & \tbd & \tbd & \tbd \\
Residual PPO & \tbd & \tbd & \tbd & \tbd & \tbd \\
\bottomrule
\end{tabular}
}
\end{table}

\subsection{Fixed-Height Balance Sweep}

% [TODO-TABLE-IV: Fixed-Height Sweep Results]
\begin{table}[ht]
\centering
\caption{Fixed-height balance sweep: performance at seven commanded heights.
Residual PPO results are mean $\pm$ std over 3 seeds.}
\label{tab:fixed_height_sweep}
\resizebox{\columnwidth}{!}{%
\begin{tabular}{cccccc}
\toprule
$h^{cmd}$ (m) & Method & Surv. (\%) & $e_h$ (m) & Pitch (deg) & Torque (Nm) \\
\midrule
\multirow{2}{*}{0.70} & LQR/IK & \tbd & \tbd & \tbd & \tbd \\
                      & Residual & \tbd & \tbd & \tbd & \tbd \\
\midrule
\multirow{2}{*}{0.65} & LQR/IK & \tbd & \tbd & \tbd & \tbd \\
                      & Residual & \tbd & \tbd & \tbd & \tbd \\
\midrule
\multirow{2}{*}{0.60} & LQR/IK & \tbd & \tbd & \tbd & \tbd \\
                      & Residual & \tbd & \tbd & \tbd & \tbd \\
\midrule
\multirow{2}{*}{0.55} & LQR/IK & \tbd & \tbd & \tbd & \tbd \\
                      & Residual & \tbd & \tbd & \tbd & \tbd \\
\midrule
\multirow{2}{*}{0.50} & LQR/IK & \tbd & \tbd & \tbd & \tbd \\
                      & Residual & \tbd & \tbd & \tbd & \tbd \\
\midrule
\multirow{2}{*}{0.45} & LQR/IK & \tbd & \tbd & \tbd & \tbd \\
                      & Residual & \tbd & \tbd & \tbd & \tbd \\
\midrule
\multirow{2}{*}{0.40} & LQR/IK & \tbd & \tbd & \tbd & \tbd \\
                      & Residual & \tbd & \tbd & \tbd & \tbd \\
\bottomrule
\end{tabular}
}
\end{table}

% [TODO-FIGURE-2: Robot posture snapshots at heights 0.40, 0.50, 0.60, 0.70 m]
\textit{[TODO-FIGURE-2: Posture snapshots across height range]}

\subsection{Commanded Height Transitions}

% [TODO-TABLE-V: Height Transition Performance]
\begin{table}[ht]
\centering
\caption{Commanded height transition performance: standing-squatting transitions.
Residual PPO results are mean $\pm$ std over 3 seeds.}
\label{tab:height_transitions}
\resizebox{\columnwidth}{!}{%
\begin{tabular}{lcccc}
\toprule
Transition & Method & Success (\%) & Settling (s) & Overshoot (m) \\
\midrule
\multirow{2}{*}{0.70 $\to$ 0.40 m} & LQR/IK & \tbd & \tbd & \tbd \\
                                    & Residual & \tbd & \tbd & \tbd \\
\midrule
\multirow{2}{*}{0.40 $\to$ 0.70 m} & LQR/IK & \tbd & \tbd & \tbd \\
                                    & Residual & \tbd & \tbd & \tbd \\
\midrule
\multirow{2}{*}{0.65 $\to$ 0.50 m} & LQR/IK & \tbd & \tbd & \tbd \\
                                    & Residual & \tbd & \tbd & \tbd \\
\bottomrule
\end{tabular}
}
\end{table}

% [TODO-FIGURE-3: Height transition time series showing commanded height,
%  actual height, pitch angle, and wheel velocity for a representative
%  transition (e.g., 0.70 → 0.40 m). Compare LQR/IK vs Residual PPO.]
\textit{[TODO-FIGURE-3: Height transition time series]}

\subsection{Push-Disturbance Recovery}

% [TODO-TABLE-VI: Push Recovery Performance]
\begin{table}[ht]
\centering
\caption{Push-disturbance recovery performance.
Residual PPO results are mean $\pm$ std over 3 seeds.}
\label{tab:push_recovery}
\resizebox{\columnwidth}{!}{%
\begin{tabular}{lccccc}
\toprule
Method & Max push (N) & Recovery (s) & Peak pitch (deg) & Surv. (\%) & Fall (\%) \\
\midrule
LQR/IK prior & \tbd & \tbd & \tbd & \tbd & \tbd \\
Residual PPO & \tbd & \tbd & \tbd & \tbd & \tbd \\
\bottomrule
\end{tabular}
}
\end{table}

% [TODO-FIGURE-4: Push recovery time series showing pitch angle, wheel velocity,
%  and residual action magnitude for a 50N push at t=3s. Compare LQR/IK vs Residual.]
\textit{[TODO-FIGURE-4: Push recovery time series]}

% [TODO-FIGURE-5: Push magnitude sweep showing survival rate vs. push force
%  (20N to 200N in 8 steps). Plot LQR/IK prior (single line) and Residual PPO
%  (mean ± std over 3 seeds).]
\textit{[TODO-FIGURE-5: Push magnitude sweep degradation curve]}

\subsection{Robustness to Model Uncertainty}

% [TODO-TABLE-VII: Robustness to Model Perturbations]
\begin{table}[ht]
\centering
\caption{Robustness to model uncertainty: friction, mass, damping, noise, and delay.
Residual PPO results are mean $\pm$ std over 3 seeds.}
\label{tab:robustness}
\begin{tabular}{lccc}
\toprule
Scenario & Method & Survival (\%) & Height RMSE (m) \\
\midrule
\multirow{2}{*}{Friction low (0.6×)} & LQR/IK & \tbd & \tbd \\
                                      & Residual & \tbd & \tbd \\
\midrule
\multirow{2}{*}{Friction high (1.4×)} & LQR/IK & \tbd & \tbd \\
                                       & Residual & \tbd & \tbd \\
\midrule
\multirow{2}{*}{Mass +20\%} & LQR/IK & \tbd & \tbd \\
                            & Residual & \tbd & \tbd \\
\midrule
\multirow{2}{*}{Damping +50\%} & LQR/IK & \tbd & \tbd \\
                               & Residual & \tbd & \tbd \\
\midrule
\multirow{2}{*}{Noise + 1-step delay} & LQR/IK & \tbd & \tbd \\
                                      & Residual & \tbd & \tbd \\
\bottomrule
\end{tabular}
\end{table}

\subsection{Residual Action Analysis}
\label{sec:residual_analysis}

% [TODO-TABLE-VIII: Residual Action Decomposition]
\begin{table}[ht]
\centering
\caption{Residual action decomposition analysis across scenarios.
Shows how the residual policy distributes corrections across joint groups.
Residual PPO results are mean $\pm$ std over 3 seeds.}
\label{tab:residual_analysis}
\resizebox{\columnwidth}{!}{%
\begin{tabular}{lcccccc}
\toprule
Scenario & Base RMS & Res. RMS & Final RMS & Res/Base & Res. rate & Sat. (\%) \\
\midrule
Nominal & \tbd & \tbd & \tbd & \tbd & \tbd & \tbd \\
Random height & \tbd & \tbd & \tbd & \tbd & \tbd & \tbd \\
Height trans. & \tbd & \tbd & \tbd & \tbd & \tbd & \tbd \\
Push recovery & \tbd & \tbd & \tbd & \tbd & \tbd & \tbd \\
Friction low & \tbd & \tbd & \tbd & \tbd & \tbd & \tbd \\
\bottomrule
\end{tabular}
}
\end{table}

% [TODO-TABLE-IX: Per-Joint-Group Residual RMS]
\begin{table}[ht]
\centering
\caption{Per-joint-group residual action RMS for random-height balance scenario.
Shows which joint groups receive the largest residual corrections.
Residual PPO results are mean $\pm$ std over 3 seeds.}
\label{tab:residual_by_joint}
\begin{tabular}{lc}
\toprule
Joint group & Residual RMS \\
\midrule
Hip roll & \tbd \\
Hip yaw & \tbd \\
Hip pitch + knee & \tbd \\
Wheels & \tbd \\
\bottomrule
\end{tabular}
\end{table}

% [TODO-FIGURE-7: Residual action distribution histogram showing the
%  distribution of residual actions across all joints and timesteps for
%  the random-height scenario. Should show that most residuals are small
%  (near zero) with occasional larger corrections.]
\textit{[TODO-FIGURE-7: Residual action distribution histogram]}

\subsection{Ablation Study}
\label{sec:ablation}

% [TODO-TABLE-X: Ablation Study Results]
\begin{table}[ht]
\centering
\caption{Ablation study: impact of key design choices.
All ablations trained for 50M steps over 3 seeds.
Baseline is the full proposed method.}
\label{tab:ablation}
\resizebox{\columnwidth}{!}{%
\begin{tabular}{lccc}
\toprule
Configuration & Survival (\%) & Height RMSE (m) & Pitch RMS (deg) \\
\midrule
\textbf{Full method (baseline)} & \tbd & \tbd & \tbd \\
\midrule
\multicolumn{4}{l}{\textit{Residual scale ablations}} \\
\quad Residual scale $\times$ 0.5 & \tbd & \tbd & \tbd \\
\quad Residual scale $\times$ 2.0 & \tbd & \tbd & \tbd \\
\quad Uniform scale (0.15 all) & \tbd & \tbd & \tbd \\
\midrule
\multicolumn{4}{l}{\textit{Observation ablations}} \\
\quad No base action in obs & \tbd & \tbd & \tbd \\
\quad Base action delayed 1 step & \tbd & \tbd & \tbd \\
\midrule
\multicolumn{4}{l}{\textit{Regularization ablations}} \\
\quad No residual penalties & \tbd & \tbd & \tbd \\
\quad 2× residual penalties & \tbd & \tbd & \tbd \\
\midrule
\multicolumn{4}{l}{\textit{Baseline comparisons}} \\
\quad LQR/IK prior only & \tbd & \tbd & \tbd \\
\quad Pure PPO (no prior) & \tbd & \tbd & \tbd \\
\bottomrule
\end{tabular}
}
\end{table}

The ablation study validates key design choices:
\begin{itemize}
    \item \textbf{Residual scale}: The per-joint residual scale vector
    (Eq.~\eqref{eq:residual_scale}) balances prior trust and correction
    authority. Halving the scale limits recovery capability; doubling it
    increases instability. Uniform scaling performs worse than joint-specific
    tuning.

    \item \textbf{Base action in observation}: Including $a_t^{\mathrm{base}}$
    in the observation (Eq.~\eqref{eq:obs_residual}) allows the residual
    policy to condition corrections on the prior's intent. Removing it
    degrades performance.

    \item \textbf{Residual regularization}: The residual-specific penalties
    (Section~\ref{sec:residual_regularization}) encourage the policy to
    rely on the prior when sufficient. Removing them increases unnecessary
    corrections; doubling them over-constrains the policy.

    \item \textbf{Pure PPO baseline}: Training a policy from scratch without
    the LQR/IK prior provides a reference for the value of structured
    initialization. [TODO: report comparative results]
\end{itemize}

% ============================================================
\section{Discussion and Limitations}
\label{sec:discussion}
% ============================================================

\subsection{Bounded Residual Framework Design}

The bounded residual approach combines structured model-based control
with learned corrections, balancing the interpretability and safety of
classical control with the adaptability of reinforcement learning.

\noindent\textbf{Structured prior as initialization.}
The LQR/IK prior provides a physically meaningful starting point for
balance and posture control. While insufficient alone (85--100\% fall
rates, Table~\ref{tab:prior_verification}), it reduces the exploration
burden on the residual policy by providing approximately correct actions
that the policy can refine rather than discover from scratch.

\noindent\textbf{Bounded corrections for safety.}
The per-joint residual scale vector (Eq.~\eqref{eq:residual_scale})
limits the magnitude of learned corrections, preventing the policy from
completely overriding the prior. This design provides a safety envelope:
even if the residual policy fails, the base controller maintains basic
balance authority.

\noindent\textbf{Base action in observation.}
Including $a_t^{\mathrm{base}}$ in the observation
(Eq.~\eqref{eq:obs_residual}) allows the residual policy to condition
its corrections on the prior's intent. This enables the policy to learn
when to trust the prior (outputting near-zero residuals) and when to
override it (outputting larger corrections).

\noindent\textbf{Progressive curriculum.}
The within-stage height curriculum stabilizes learning by starting near
the nominal standing pose and expanding toward more challenging
configurations. The height-consistent natural pose reward constrains
the joint-space solution to physically meaningful postures at each
target height.

\noindent\textbf{Shared low-level controller.}
The PID low-level controller decouples the residual policy from
actuator-level dynamics and is shared with the LQR/IK prior, ensuring
that observed performance differences reflect controller intelligence
rather than interface differences.

\subsection{Sim-to-Real Considerations}
\label{sec:sim2real}
Although the present work is conducted entirely in simulation,
the framework is designed with hardware transfer in mind.

\noindent\textbf{Action representation.}
The policy outputs PD position targets rather than raw torques.
This action representation has been shown to transfer reliably across
the sim-to-real gap for legged locomotion~\cite{hwangbo2019learning,rudin2022walk},
as it implicitly low-pass filters the control signal and decouples the policy
from precise actuator gain tuning.

\noindent\textbf{Domain randomization coverage.}
The randomization ranges in Table~\ref{tab:domain_rand} were chosen to
separate nominal balance learning from later robustness fine-tuning.
The \texttt{balance} stage uses moderate perturbations to avoid learning
spurious wheel-spin or drift strategies before a clean stationary posture is
established.
The \texttt{balance\_robust} stage then widens friction, mass, and damping
variation to cover floor surfaces from polished tile to rubber mat, battery
charge variation and light payloads, actuator damping uncertainty, and
communication latency through optional action delay.

\noindent\textbf{Observation availability.}
The policy uses body-frame linear velocity, which requires a velocity
estimator on hardware.
Should a reliable estimator be unavailable, the velocity channel can
be dropped, reducing the observation to 39 dimensions.
The policy can be retrained under this reduced observation for
direct deployment, at a modest cost in height-tracking precision.

\noindent\textbf{Transfer plan.}
Planned hardware deployment will proceed in three stages:
(i)~system identification of actuator friction and rotor inertia
to refine the simulator;
(ii)~policy deployment at fixed nominal height ($h^{cmd}=0.65$\,m)
with a fall-safe watchdog;
(iii)~progressive extension to variable-height commands
following the same Phase~A/B/C ordering as training.

\subsection{Current Limitations}
\begin{enumerate}
    \item \textbf{No quantitative results yet.}
    Residual policy training has not been completed.
    All result sections contain TODO placeholders.
    The LQR/IK prior has been evaluated (Table~\ref{tab:prior_verification}),
    confirming its limitations and motivating the residual approach,
    but the residual policy training and evaluation remain future work.

    \item \textbf{Multi-seed training planned, not yet complete.}
    Final metrics will aggregate runs with seeds 42, 113, and 999.
    Quantitative results will report mean $\pm$ std across all three seeds.

    \item \textbf{No sim-to-real transfer (yet).}
    All work is simulation-only.
    Although the model matches a physical platform, hardware transfer
    has not been attempted.
    Sensor-noise and action-delay settings follow established
    sim-to-real protocols~\cite{tobin2017domain,peng2018simtoreal,tan2018simtoreal},
    but validation remains future work.

    \item \textbf{Limited prior validation.}
    The LQR/IK prior has been evaluated at fixed heights, showing
    85--100\% fall rates. However, the prior has not been validated
    under push disturbances, friction variations, or commanded height
    transitions. These evaluations will be included in the final
    experimental suite.

    \item \textbf{Stand-up recovery not evaluated.}
    Recovery from fallen or non-nominal initial poses is mentioned
    as a potential extension but has not been implemented or evaluated.
    All claims about stand-up recovery are structural, not
    performance-validated.

    \item \textbf{Pure PPO baseline not yet trained.}
    The ablation study includes a pure PPO baseline (no LQR/IK prior)
    for comparison, but this baseline has not been trained.
    Comparative claims about the value of the structured prior
    await experimental validation.
    % [TODO-CLAIM: only state LQR inferiority at large disturbances after
    %  push_recovery and full_range scenarios are evaluated]

    \item \textbf{Conservative curriculum advancement.}
    The within-stage height curriculum evaluates roughly every
    $10^7$ environment steps using 256 fixed-horizon greedy episodes.
    Advancement additionally requires training reward, evaluation return,
    success rate, and fall-rate gates to pass.
    This reduces premature range expansion but may slow progress when a
    policy is close to satisfying the next curriculum level.
\end{enumerate}

% ============================================================
\section{Conclusion}
% ============================================================
This paper presents a bounded residual reinforcement learning framework
for height-adaptive balance control of a 10-actuated wheeled biped robot.
The approach combines a height-dependent LQR/IK prior with a learned
residual policy that outputs bounded corrections, balancing the
interpretability and safety of classical control with the adaptability
of reinforcement learning.

The core technical contributions are:
(i)~a height-dependent LQR/IK prior that provides structured balance
and posture commands, validated to have 85--100\% fall rates across
the commanded height range $[0.40, 0.70]\,\mathrm{m}$, motivating the
need for learned corrections;
(ii)~a bounded residual policy trained via PPO that observes the
prior's action and outputs per-joint-scaled corrections, enabling
the policy to learn when to trust the prior and when to override it;
(iii)~a 52-dimensional residual observation space that includes the
base action, allowing the policy to condition corrections on the
prior's intent;
(iv)~residual-specific reward regularization that encourages the
policy to rely on the prior when sufficient;
(v)~a within-stage height curriculum that expands the commanded
height range across 29 levels, gated by greedy evaluation;
(vi)~a height-consistent natural pose reward that maps commanded
height to target joint angles, preventing degenerate balancing
strategies.

The framework is designed for sim-to-real transfer: PD position targets
as the action representation, comprehensive domain randomization covering
friction, mass, joint damping, and sensor noise, and an observation
space that degrades gracefully when linear velocity is unavailable.

% [TODO-RESULT: Once training is complete, replace the following sentence
%  with a summary of key quantitative findings, e.g.:
%  "Experimental results show that the residual policy achieves X\% survival
%   rate across the full [0.40, 0.70] m range and recovers from pushes up to Y N,
%   outperforming the LQR/IK prior on all height-adaptive and disturbance metrics.
%   Ablation studies confirm the value of including the base action in the
%   observation and the per-joint residual scale design."]
Experimental evaluation and quantitative comparison with the LQR/IK prior
are ongoing; results will be reported after residual policy training
completes over three independent seeds (42, 113, 999).

Future work will focus on hardware deployment via the staged transfer
plan described in Section~\ref{sec:sim2real}, training and evaluation
of the pure PPO baseline for comparative analysis, and extension to
stand-up recovery and robust push-disturbance scenarios.

\section*{Acknowledgment}
% [TODO: Add acknowledgments if applicable.]

\bibliographystyle{IEEEtran}
\begin{thebibliography}{99}

\bibitem{ppo}
J. Schulman, F. Wolski, P. Dhariwal, A. Radford, and O. Klimov,
``Proximal Policy Optimization Algorithms,''
\emph{arXiv preprint arXiv:1707.06347}, 2017.

\bibitem{mujoco}
E. Todorov, T. Erez, and Y. Tassa,
``MuJoCo: A physics engine for model-based control,''
in \emph{Proc. IEEE/RSJ Int. Conf. Intelligent Robots and Systems (IROS)},
2012, pp. 5026--5033.

\bibitem{mujoco_mjx_docs}
``MuJoCo Documentation: MuJoCo XLA (MJX),''
[Online]. Available:
\url{https://mujoco.readthedocs.io/en/latest/mjx.html}.
Accessed: Apr. 2026.

\bibitem{jax}
J. Bradbury \emph{et al.},
``JAX: composable transformations of Python+NumPy programs,''
2018. [Online]. Available: \url{http://github.com/jax-ml/jax}

\bibitem{flax}
J. Heek, A. Levskaya, A. Oliver, \emph{et al.},
``Flax: A neural network library and ecosystem for JAX,''
2024. [Online]. Available: \url{http://github.com/google/flax}

\bibitem{gae}
J. Schulman, P. Moritz, S. Levine, M. Jordan, and P. Abbeel,
``High-Dimensional Continuous Control Using Generalized Advantage
Estimation,''
in \emph{Proc. Int. Conf. Learning Representations (ICLR)}, 2016.

\bibitem{welford1962}
B. P. Welford,
``Note on a Method for Calculating Corrected Sums of Squares
and Products,''
\emph{Technometrics}, vol.~4, no.~3, pp. 419--420, 1962.

\bibitem{tobin2017domain}
J. Tobin, R. Fong, A. Ray, J. Schneider, W. Zaremba, and P. Abbeel,
``Domain Randomization for Transferring Deep Neural Networks from
Simulation to the Real World,''
in \emph{Proc. IEEE/RSJ Int. Conf. Intelligent Robots and Systems (IROS)},
2017, pp. 23--30.

\bibitem{bengio2009curriculum}
Y. Bengio, J. Louradour, R. Collobert, and J. Weston,
``Curriculum Learning,''
in \emph{Proc. Int. Conf. Machine Learning (ICML)}, 2009, pp. 41--48.

\bibitem{klemm2019ascento}
V. Klemm \emph{et al.},
``Ascento: A Two-Wheeled Jumping Robot,''
in \emph{Proc. IEEE Int. Conf. Robotics and Automation (ICRA)}, 2019.

\bibitem{klemm2020lqrwbc}
V. Klemm, A. Morra, L. Gulich, \emph{et al.},
``LQR-Assisted Whole-Body Control of a Wheeled Bipedal Robot
with Kinematic Loops,''
\emph{IEEE Robotics and Automation Letters},
vol.~5, no.~2, pp. 3745--3752, 2020.

\bibitem{hsu2022wbrfuzzy}
C.-F. Hsu, B.-R. Chen, and Z.-L. Lin,
``Implementation and Control of a Wheeled Bipedal Robot Using
a Fuzzy Logic Approach,''
\emph{Actuators}, vol.~11, no.~12, p. 357, 2022.

\bibitem{li2021wheelbiped}
X. Li \emph{et al.},
``Dynamic wheeled motion control of wheel-biped transformable robots,''
\emph{Biomimetic Intelligence and Robotics}, 2021.

\bibitem{rudin2022walk}
N. Rudin, D. Hoeller, P. Reist, and M. Hutter,
``Learning to Walk in Minutes Using Massively Parallel Deep Reinforcement Learning,''
in \emph{Proc. Conf. Robot Learning (CoRL)}, vol.~164, pp. 91--100, 2022.

\bibitem{lee2020learning}
J. Lee, J. Hwangbo, L. Wellhausen, V. Koltun, and M. Hutter,
``Learning Quadrupedal Locomotion over Challenging Terrain,''
\emph{Science Robotics}, vol.~5, no.~47, p.~eabc5986, 2020.

\bibitem{kumar2021rma}
A. Kumar, Z. Fu, D. Pathak, and J. Malik,
``RMA: Rapid Motor Adaptation for Legged Robots,''
in \emph{Proc. Robotics: Science and Systems (RSS)}, 2021.

\bibitem{hwangbo2019learning}
J. Hwangbo \emph{et al.},
``Learning Agile and Dynamic Motor Skills for Legged Robots,''
\emph{Science Robotics}, vol.~4, no.~26, p.~eaau5872, 2019.

\bibitem{miki2022wild}
T. Miki \emph{et al.},
``Learning Robust Perceptive Locomotion for Quadrupedal Robots in the Wild,''
\emph{Science Robotics}, vol.~7, no.~62, p.~eabk2822, 2022.

\bibitem{lee2024robustnavigation}
J. Lee \emph{et al.},
``Learning Robust Autonomous Navigation and Locomotion for Wheeled-Legged Robots,''
\emph{Science Robotics}, vol.~9, 2024.

\bibitem{vollenweider2023advancedskills}
E. Vollenweider \emph{et al.},
``Advanced Skills through Multiple Adversarial Motion Priors in Reinforcement Learning,''
in \emph{Proc. IEEE Int. Conf. Robotics and Automation (ICRA)}, pp. 5120--5126, 2023.

\bibitem{chamorro2024stairclimbing}
I. Chamorro, J. Arrizabalaga, and M. Hutter,
``Reinforcement Learning for Blind Stair Climbing with Legged and Wheeled-Legged Robots,''
in \emph{Proc. IEEE Int. Conf. Robotics and Automation (ICRA)}, 2024.

\bibitem{whleaper2024}
J. Zhu, S. Zhao, and E. Yoshida,
``Whleaper: A 10-DOF Flexible Bipedal Wheeled Robot,''
in \emph{Proc. IEEE/RSJ Int. Conf. Intelligent Robots and Systems (IROS)},
pp. 11272--11277, 2024.

\bibitem{zhu2022hierarchical}
J. Zhu and M. Hayashibe,
``Reinforcement Learning Based Hierarchical Control for Path Tracking of a
Wheeled Bipedal Robot with Sim-to-Real Framework,''
in \emph{Proc. IEEE/SICE Int. Symp. System Integration (SII)}, 2022.

\bibitem{diablo2024}
S. Liu \emph{et al.},
``DIABLO: A 6-DoF Wheeled Bipedal Robot Composed Entirely of Direct-Drive Joints,''
\emph{arXiv preprint arXiv:2407.21500}, 2024.

\bibitem{portelas2020acl}
R. Portelas, C. Colas, L. Weng, K. Hofmann, and P.-Y. Oudeyer,
``Automatic Curriculum Learning for Deep RL: A Short Survey,''
in \emph{Proc. Int. Joint Conf. Artificial Intelligence (IJCAI)}, 2020.

\bibitem{portelas2020teacher}
R. Portelas, C. Colas, K. Hofmann, and P.-Y. Oudeyer,
``Teacher Algorithms for Curriculum Learning of Deep RL in Continuously
Parameterized Environments,''
in \emph{Proc. Conf. Robot Learning (CoRL)}, vol.~100, pp. 835--853, 2020.

\bibitem{wang2019poet}
R. Wang, J. Lehman, J. Clune, and K. O. Stanley,
``Paired Open-Ended Trailblazer (POET): Endlessly Generating Increasingly
Complex and Diverse Learning Environments and Their Solutions,''
\emph{arXiv preprint arXiv:1901.01753}, 2019.

\bibitem{peng2018simtoreal}
X. B. Peng, M. Andrychowicz, W. Zaremba, and P. Abbeel,
``Sim-to-Real Transfer of Robotic Control with Dynamics Randomization,''
in \emph{Proc. IEEE Int. Conf. Robotics and Automation (ICRA)},
pp. 3803--3810, 2018.

\bibitem{tan2018simtoreal}
J. Tan \emph{et al.},
``Sim-to-Real: Learning Agile Locomotion for Quadruped Robots,''
in \emph{Proc. Robotics: Science and Systems (RSS)}, 2018.

\end{thebibliography}

\end{document}
